\documentclass[11pt]{article}
\usepackage[margin=0.9in]{geometry}
\usepackage{graphicx}
\usepackage{booktabs}
\usepackage{longtable}
\usepackage{hyperref}
\usepackage{array}
\title{CNN-SAE analysis report}
\author{Codex playbook analysis pass}
\date{2026-06-08}
\begin{document}
\maketitle
\tableofcontents
\section{Executive summary}
This analysis pass treats code and run artifacts as primary evidence and writes all new outputs under analysis/. The strongest artifact-backed positives are: the FieldSAE winner remains highly compressive, the Q3 retained-fraction bottleneck persists in the pareto sweep, full-BPTT re-grounded chaining reaches 0.7010 against a 0.7122 hybrid bound, raw-only rollout-aligned retraining reaches 0.6821, the original ViT transfer artifact attains 0.9440 prediction agreement, the completed 75-run ViT-small sweep recovers stable winner settings with mean prediction agreement up to 0.9307, and the completed ViT-base sweep reaches a much stronger stable regime with mean prediction agreement up to 0.9851 at block 10 and 12 percent retained.


The most important caveat is a confirmed repeated-test-exposure pattern across the CIFAR training families: core training scripts evaluate on the test split every epoch. That does not undercut the central mechanistic claim, because this project uses that split as an analysis corpus for probing how the frozen CNN executes. It does mean the reported CIFAR top-1 values are corpus-conditioned faithfulness measurements rather than untouched external evaluation numbers.


\section{Experiment catalog}
\begin{longtable}{p{0.21\linewidth}p{0.15\linewidth}p{0.18\linewidth}p{0.36\linewidth}}
\toprule
Experiment & Stage & Parallel group & Purpose \\
\midrule
\endhead
phase1\_anchor & Phase 1 & section x retained\_fraction & Anchor reconstruction with VectorSAE plus global masking. \\
phase1\_random\_baseline & Phase 1 & selected section x retained\_fraction & Random-dictionary sanity checks against the trained SAE anchor. \\
phase1\_pca\_baseline & Phase 1 & section x PCA rank & Linear PCA baseline at matched or larger coefficient budgets. \\
phase1\_pareto & Phase 1 addendum & section x retained\_fraction & Fine retained-fraction sweep on the FieldSAE winner. \\
phase2\_grid & Phase 2 & cell x section x mask\_budget & Branch grid over Field/Vector SAE and global/RF masking. \\
phase3\_transitions & Phase 3 & adjacent section transition & Family I adjacent sparse-code transitions and frozen-block hybrid bound. \\
phase3\_q3q4\_strong\_retry & Phase 3 addendum & single focused retry & Capacity-scaled retry of the Q3 to Q4 transition bottleneck. \\
phase3b\_chain & Phase 3b & epoch progression & Chain-aware scheduled-sampling training with re-grounding. \\
phase3b\_bptt & Phase 3b addendum & epoch progression & Full-BPTT chain-aware training through differentiable re-grounding. \\
phase2\_rf\_fair\_changed & Phase 3b changed & section x retained\_fraction & Fair-budget RF-local masking rerun. \\
phase3b\_raw\_only\_changed & Phase 3b changed & epoch progression & Rollout-aligned raw-carrier chain retraining without re-grounding. \\
phase1\_patch\_overlap\_changed & Phase 3b changed & section x patch\_multiplier x retained\_fraction & Overlapping patchwise SAE retry. \\
phase1\_patch\_depth\_changed & Phase 3b changed & section x depth\_schedule x retained\_fraction & Depth-aware overlapping patchwise SAE sweep. \\
phase1\_patch\_disjoint\_changed & Phase 3b changed & section x patch\_multiplier & Disjoint patchwise SAE sweep. \\
phase4\_r20c10 & Phase 4 & section x retained\_fraction & ResNet-20 on CIFAR-10 transfer grid for the winner recipe. \\
phase4\_seeds & Phase 4 & section x seed & Seed-stability grid on the ResNet-56 / CIFAR-100 winner setup. \\
phase4\_taxonomy & Phase 4 & representation condition & Transition taxonomy comparing sparse, dense, and cross carriers. \\
phase4\_r110 & Phase 4 & section & ResNet-110 reconstruction sweep on CIFAR-100. \\
phase4\_vit\_transfer & Phase 4 & single ViT transfer artifact & FieldSAE transfer to a pretrained ViT block on Imagenette. \\
phase4\_vit\_sweep & Phase 4 addendum & block x retained\_fraction x seed & Completed ViT-small sweep on Imagenette over 75 cells. \\
phase4\_vit\_base & Phase 4 addendum & block x retained\_fraction x seed & Completed ViT-base sweep on Imagenette over 60 cells. \\
tier1\_intervention & Tier 1 follow-up & adjacent section transition & Parent-to-child intervention audit of importance and locality. \\
tier2\_feature\_audit & Tier 2 follow-up & section & Dictionary dead-feature and duplicate-feature audit. \\
\bottomrule
\end{longtable}
\section{Per-experiment notes}
\subsection{phase1\_anchor}
Purpose and hypothesis: A sparse transform-domain code can reconstruct each section and preserve downstream behavior.


Designed method: VectorSAE with global TopK masking over cached ResNet-56 / CIFAR-100 activations.


Designed procedure: Train one SAE per section over fractions 2, 5, and 10 percent.


Declared success criteria: Beat random and PCA baselines while staying within roughly 1 percentage point of original top-1.


Outcome classification: Success (as hoped). Strong reconstruction and downstream preservation across sections. The CIFAR top-1 values should be read as corpus-conditioned faithfulness measurements rather than untouched external estimates.


\subsection{phase1\_random\_baseline}
Purpose and hypothesis: The anchor result should not be reproducible with an untrained dictionary.


Designed method: Freeze a random SAE dictionary and evaluate the same splice metrics.


Designed procedure: Run selected random-baseline cells on the anchor sections.


Declared success criteria: Random baseline should collapse toward chance.


Outcome classification: Success (as hoped). Random dictionary runs collapse toward chance, matching the intended sanity check.


\subsection{phase1\_pca\_baseline}
Purpose and hypothesis: A learned sparse transform should outperform linear PCA at similar budget.


Designed method: Per-section PCA over normalized cached activations.


Designed procedure: Evaluate several ranks on the cached test activations.


Declared success criteria: PCA should underperform the learned SAE at matched-or-larger budget.


Outcome classification: Success (as hoped). PCA trails the learned SAE across the reported sections.


\subsection{phase1\_pareto}
Purpose and hypothesis: The FieldSAE winner has a clear retained-fraction knee and a persistent Q3 bottleneck.


Designed method: FieldSAE with global masking over fractions 1, 2, 3, 5, 8, 12, and 20 percent.


Designed procedure: Sweep all sections on the same backbone with one final metric per cell.


Declared success criteria: Recover monotone sparsity-fidelity curves and identify the operational sweet spot.


Outcome classification: Success (as hoped). Recovered monotone sparsity-fidelity curves with Q3 as the persistent late-budget bottleneck.


\subsection{phase2\_grid}
Purpose and hypothesis: An expressive FieldSAE should beat the VectorSAE, especially at aggressive sparsity and under RF-local masking.


Designed method: Grid over Field versus Vector and Global versus RF masking.


Designed procedure: Evaluate all planned section and budget combinations for the three surviving cells.


Declared success criteria: Identify a clear winner architecture for the hierarchy stages.


Outcome classification: Success (as hoped). FieldSAE clearly dominates the VectorSAE on the aggressive-sparsity and RF-local branches.


\subsection{phase3\_transitions}
Purpose and hypothesis: Sparse code at one section can generate the next section, and the frozen-block hybrid should upper-bound learned transitions.


Designed method: Train adjacent transition predictors between independent SAEs.


Designed procedure: Train four adjacent predictors and compare learned splice accuracy against the hybrid bound.


Declared success criteria: Strong single-step transitions and an informative hybrid-vs-learned gap.


Outcome classification: Success (as hoped). Single-step transitions work, and the hybrid upper bound shows strong causal sufficiency. The CIFAR top-1 values should be read as corpus-conditioned faithfulness measurements rather than untouched external estimates.


\subsection{phase3\_q3q4\_strong\_retry}
Purpose and hypothesis: The hard Q3 to Q4 step may yield to a larger predictor.


Designed method: Retry the Q3 to Q4 transition with a deeper and wider predictor.


Designed procedure: Train one focused retry for the bottleneck edge.


Declared success criteria: Close the gap between learned and hybrid Q3 to Q4 performance.


Outcome classification: Undesirable result. The larger predictor makes the Q3 to Q4 edge worse, not better.


\subsection{phase3b\_chain}
Purpose and hypothesis: Scheduled sampling on re-grounded predicted inputs should stabilize the learned chain.


Designed method: Chain-aware training over the four adjacent predictors.


Designed procedure: Warm-start from Phase 3 and track chain metrics through training.


Declared success criteria: Lift the re-grounded chain above the Family I baseline.


Outcome classification: Success (as hoped). Chain-aware training lifts the re-grounded chain well above the Family I baseline.


\subsection{phase3b\_bptt}
Purpose and hypothesis: Training through differentiable re-grounding should narrow the gap to the hybrid upper bound.


Designed method: Full BPTT through the chain with in-graph re-grounding.


Designed procedure: Repeat the chain-aware training with --bptt enabled.


Declared success criteria: Push the re-grounded chain closer to the hybrid chain than the original Phase 3b run.


Outcome classification: Success (as hoped). Full BPTT closes most of the gap to the hybrid chain.


\subsection{phase2\_rf\_fair\_changed}
Purpose and hypothesis: The earlier RF-local story may have been distorted by unfair budgeting.


Designed method: Changed trainer with fair per-window RF fraction accounting.


Designed procedure: Run all sections at 2, 5, and 10 percent.


Declared success criteria: Show whether RF-locality remains viable under fair budgeting.


Outcome classification: Success (unexpected). Fair RF budgeting rescues a branch that earlier looked much weaker.


\subsection{phase3b\_raw\_only\_changed}
Purpose and hypothesis: Raw chaining may be recoverable if the model is trained on its own rollout distribution.


Designed method: Raw-carrier chain training with scheduled sampling, BPTT, and clipping.


Designed procedure: Train one chain variant without re-grounding in the carrier path.


Declared success criteria: Move the raw chain well above chance.


Outcome classification: Success (unexpected). Rollout-aligned retraining recovers a strong raw chain instead of chance collapse.


\subsection{phase1\_patch\_overlap\_changed}
Purpose and hypothesis: The old VectorSAE baseline failed partly because it lacked local spatial context.


Designed method: Patchwise overlapping sparse coding with several patch multipliers.


Designed procedure: Sweep sections, fractions, and patch multipliers.


Declared success criteria: Recover stronger low-budget performance than the pointwise vector baseline.


Outcome classification: Success (unexpected). Patchwise overlap materially improves over the pointwise vector baseline.


\subsection{phase1\_patch\_depth\_changed}
Purpose and hypothesis: Patch size should shrink with depth instead of staying uniformly large.


Designed method: Depth-aware overlapping patch schedules.


Designed procedure: Sweep schedules a-d over sections and fractions.


Declared success criteria: Find a better schedule than uniform large patches.


Outcome classification: Success (unexpected). Depth-aware schedules complete cleanly and broadly support the small-context story.


\subsection{phase1\_patch\_disjoint\_changed}
Purpose and hypothesis: Disjoint tilings may behave differently from overlapping patch scans.


Designed method: Disjoint patchwise sparse coding at a single retained fraction over several patch scales.


Designed procedure: Run all sections at 5 percent for four patch multipliers.


Declared success criteria: Determine whether non-overlapping patches are competitive with overlap.


Outcome classification: Undesirable result. Coarse disjoint tilings fail badly; only the smallest disjoint scale remains competitive.


\subsection{phase4\_r20c10}
Purpose and hypothesis: The winner recipe should survive a smaller backbone and dataset shift.


Designed method: FieldSAE with global masking on ResNet-20 / CIFAR-10.


Designed procedure: Run all sections at 2, 5, and 10 percent.


Declared success criteria: Preserve near-original top-1 across sections on the new backbone/dataset.


Outcome classification: Success (as hoped). The winner recipe survives the ResNet-20 / CIFAR-10 transfer grid.


\subsection{phase4\_seeds}
Purpose and hypothesis: The winner story should not be a single-seed accident.


Designed method: Repeat the winner SAE setup across seeds 0, 1, and 2.


Designed procedure: Train all five sections at 5 percent across the three seeds.


Declared success criteria: Low enough variance that the qualitative story survives.


Outcome classification: Success (as hoped). Seed variation is small enough that the qualitative story survives.


\subsection{phase4\_taxonomy}
Purpose and hypothesis: Sparse code is the strongest transition carrier among sparse, dense, and cross alternatives.


Designed method: Controlled Q4 to Q5 comparison across three representation families.


Designed procedure: Measure downstream preservation, prediction agreement, and KL for each condition.


Declared success criteria: Sparse should dominate the alternative carriers.


Outcome classification: Success (as hoped). Sparse is the best transition carrier among the tested alternatives.


\subsection{phase4\_r110}
Purpose and hypothesis: The winner recipe should transfer to a deeper CNN.


Designed method: FieldSAE global masking on ResNet-110 / CIFAR-100 at 5 percent.


Designed procedure: Run all five sections once.


Declared success criteria: Preserve near-original accuracy and recover the same bottleneck pattern.


Outcome classification: Success (as hoped). The winner recipe transfers to a deeper CNN and preserves the Q3 bottleneck pattern.


\subsection{phase4\_vit\_transfer}
Purpose and hypothesis: The FieldSAE path should transfer to a ViT token grid.


Designed method: FieldSAE over block-level ViT patch tokens on Imagenette.


Designed procedure: Train on Imagenette train split and evaluate self-consistency on val.


Declared success criteria: High prediction agreement and low KL at sparse retained fraction.


Outcome classification: Success (as hoped). The ViT transfer artifact achieves high prediction agreement with low KL at 5 percent retained.


\subsection{phase4\_vit\_sweep}
Purpose and hypothesis: The ViT transfer result should remain stable across multiple blocks, sparsity budgets, and seeds rather than depending on one lucky cell.


Designed method: FieldSAE over ViT-small patch-token grids on Imagenette, sweeping blocks 2, 4, 6, 8, and 10; retained fractions 1, 2, 5, 8, and 12 percent; and seeds 0, 1, and 2.


Designed procedure: Run 75 total cells and summarize mean and seed spread per block/fraction pair.


Declared success criteria: Recover a broad region of high prediction agreement with low KL, and identify stable winner settings rather than a single one-off artifact.


Outcome classification: Success (as hoped). All 75 cells complete successfully, the mean prediction agreement across the sweep is 0.8914, and the best stable setting reaches 0.9307 mean prediction agreement with 0.0034 seed std and 0.0361 mean KL.


\subsection{phase4\_vit\_base}
Purpose and hypothesis: The clean ViT transfer path should scale to a larger pretrained model while preserving a stable sparse reconstruction regime.


Designed method: FieldSAE over ViT-base patch-token grids on Imagenette, sweeping blocks 2, 4, 6, 8, and 10; requested retained fractions 2, 5, 8, and 12 percent; and seeds 0, 1, and 2.


Designed procedure: Run 60 total cells and summarize both stable winner settings and any seed-fragile failure regions.


Declared success criteria: Recover a larger-model regime with near-perfect prediction agreement while keeping the sweep interpretable rather than hiding instability inside a single global average.


Outcome classification: Success (unexpected). The larger model supports an extremely strong stable late-block regime, with the best stable setting reaching 0.9851 mean prediction agreement and 0.0012 seed std, but many mid-block cells become sharply bimodal across seeds and several collapse while retaining far less than their requested sparsity budget.


\subsection{tier1\_intervention}
Purpose and hypothesis: Sparse parents should be causally important and spatially localized for their children.


Designed method: Importance and locality intervention metrics across adjacent transitions.


Designed procedure: Ablate top versus random parents and compare localized child responses.


Declared success criteria: Large importance ratios and strong locality concentration.


Outcome classification: Success (as hoped). Interventions show strong importance and locality concentration, though they inherit upstream SAE validity risks.


\subsection{tier2\_feature\_audit}
Purpose and hypothesis: The learned dictionary should be healthy rather than dead or duplicated.


Designed method: Audit dead features, duplicate features, and mean atom cosine.


Designed procedure: Run the audit on representative mid and deep sections.


Declared success criteria: Low dead-feature and duplicate-feature fractions.


Outcome classification: Success (as hoped). Dead-feature and duplicate-feature rates are low in the audited sections.


\section{Cross-experiment analysis}
Across the reconstruction families, the clearest retained-fraction knee appears in the FieldSAE pareto sweep. Q3 reaches 0.7186 at 12 percent retained, while most other sections flatten much earlier.


Across the chain families, re-grounding remains the strongest general mechanism. The scheduled-sampling Phase 3b chain reaches 0.6624, the BPTT variant reaches 0.7010, and the raw-only changed run reaches 0.6821 on its raw carrier path.


Phase 4 strengthens the transfer story materially. The cleanest surviving split discipline is the ViT family, where the original artifact reaches 0.9440 prediction agreement on the Imagenette val split, the completed 75-run ViT-small sweep shows a broad stable regime with best mean prediction agreement 0.9307 at block 10 and 12 percent retained, and the completed ViT-base sweep shows that a larger model can reach 0.9851 mean prediction agreement in its stable regime while also exposing substantial off-regime seed fragility.


\begin{figure}[ht]
\centering
\includegraphics[width=0.95\linewidth]{figures/figure\_phase1\_pareto.pdf}
\caption{Phase 1 reconstruction curves. The anchor VectorSAE needs more budget on early sections, while the FieldSAE pareto sweep reaches its useful regime by roughly 3 to 8 percent on most sections. Interpret CIFAR top-1 values with caution because the training scripts repeatedly evaluate on the test split.}
\end{figure}
\begin{figure}[ht]
\centering
\includegraphics[width=0.95\linewidth]{figures/figure\_phase2\_grid.pdf}
\caption{Phase 2 branch grid. FieldSAE dominates the aggressive-sparsity branch, and RF-local masking is only viable with the expressive FieldSAE. CIFAR top-1 values here are corpus-conditioned faithfulness measurements because this evaluation corpus was repeatedly probed during training.}
\end{figure}
\begin{figure}[ht]
\centering
\includegraphics[width=0.95\linewidth]{figures/figure\_phase3\_chain\_and\_edges.pdf}
\caption{Hierarchy-focused figures. Re-grounded chain training improves steadily, especially with full BPTT, while Q3 to Q4 remains the weakest single-step transition. CIFAR top-1 values here are corpus-conditioned faithfulness measurements because this evaluation corpus was repeatedly probed during training.}
\end{figure}
\begin{figure}[ht]
\centering
\includegraphics[width=0.95\linewidth]{figures/figure\_changed\_patch\_suites.pdf}
\caption{Changed-suite patch sweeps. Overlap helps, depth-aware schedules are broadly consistent with the small-context story, and coarse disjoint tilings fail badly. CIFAR top-1 values here are corpus-conditioned faithfulness measurements because this evaluation corpus was repeatedly probed during training.}
\end{figure}
\begin{figure}[ht]
\centering
\includegraphics[width=0.95\linewidth]{figures/figure\_phase4\_extensions.pdf}
\caption{Phase 4 extensions. The CIFAR transfer and seed figures strengthen the winner-path story but should be read as corpus-conditioned faithfulness measurements because this evaluation corpus was repeatedly probed during training. The ViT prediction-agreement marker comes from the cleanest surviving split discipline in the repo.}
\end{figure}
\begin{figure}[ht]
\centering
\includegraphics[width=0.95\linewidth]{figures/fig14\_vit\_sweep\_heatmaps.pdf}
\caption{Completed ViT-small sweep heatmaps over block and retained fraction. Mean prediction agreement, mean KL, seed-level agreement spread, and mean relative reconstruction error show a broad stable regime rather than a single lucky transfer artifact.}
\end{figure}
\begin{figure}[ht]
\centering
\includegraphics[width=0.95\linewidth]{figures/fig15\_vit\_sweep\_curves.pdf}
\caption{Completed ViT-small sweep curves. Mean metrics versus retained fraction show nearly monotone improvement with larger budgets and reveal that blocks 4 and 10 are the strongest aggregate regions in this sweep.}
\end{figure}
\begin{figure}[ht]
\centering
\includegraphics[width=0.95\linewidth]{figures/fig16\_vit\_sweep\_pareto.pdf}
\caption{Completed ViT-small sweep Pareto views. The strongest settings jointly achieve high prediction agreement and low KL / relative reconstruction error rather than trading one metric off catastrophically against another.}
\end{figure}
\begin{figure}[ht]
\centering
\includegraphics[width=0.95\linewidth]{figures/fig20\_vit\_base\_heatmaps.pdf}
\caption{Completed ViT-base sweep heatmaps. The larger model has an exceptionally strong stable regime at block 10, but several mid-block settings show high seed spread or outright collapse counts, so mean performance alone is not a sufficient summary.}
\end{figure}
\begin{figure}[ht]
\centering
\includegraphics[width=0.95\linewidth]{figures/fig20\_vit\_base\_seed\_traces.pdf}
\caption{Completed ViT-base sweep seed traces. Stable late-block settings stay tightly clustered across seeds, while multiple other settings bifurcate sharply into near-perfect and near-zero prediction-agreement outcomes.}
\end{figure}
\begin{figure}[ht]
\centering
\includegraphics[width=0.95\linewidth]{figures/fig20\_vit\_base\_fraction\_drift.pdf}
\caption{Completed ViT-base requested-versus-realized sparsity diagnostics. Several collapse-prone seeds retain far less than the requested fraction, so the analysis groups settings by the requested fraction recovered from the run directory names rather than by rounded realized sparsity alone.}
\end{figure}
\begin{figure}[ht]
\centering
\includegraphics[width=0.95\linewidth]{figures/figure\_intervention\_and\_audit.pdf}
\caption{Tier follow-up audits. The intervention ratios are large and the audited dictionaries have low dead-feature and duplicate-feature rates. These analyses inherit upstream CIFAR training validity risks because they depend on previously trained SAEs.}
\end{figure}
\section{Flags and risks}
\begin{itemize}
\item Material deviation: This is a confirmed repeated-test-exposure pattern. It does not undercut the mechanistic objective of matching the frozen CNN on the chosen analysis corpus, but it does mean the resulting fidelity numbers are corpus-conditioned rather than untouched external measurements.
\item Benign deviation: This is a documentation lag rather than a code-path mismatch.
\item Benign deviation: Another documentation lag, but important because it changes the patchwise conclusion.
\item Benign deviation: The artifact parses cleanly, but the launcher metadata is inconsistent.
\end{itemize}
\section{What worked / what we hoped for}
\begin{itemize}
\item phase1\_anchor: Strong reconstruction and downstream preservation across sections. The CIFAR top-1 values should be read as corpus-conditioned faithfulness measurements rather than untouched external estimates.
\item phase1\_random\_baseline: Random dictionary runs collapse toward chance, matching the intended sanity check.
\item phase1\_pca\_baseline: PCA trails the learned SAE across the reported sections.
\item phase1\_pareto: Recovered monotone sparsity-fidelity curves with Q3 as the persistent late-budget bottleneck.
\item phase2\_grid: FieldSAE clearly dominates the VectorSAE on the aggressive-sparsity and RF-local branches.
\item phase3\_transitions: Single-step transitions work, and the hybrid upper bound shows strong causal sufficiency. The CIFAR top-1 values should be read as corpus-conditioned faithfulness measurements rather than untouched external estimates.
\item phase3b\_chain: Chain-aware training lifts the re-grounded chain well above the Family I baseline.
\item phase3b\_bptt: Full BPTT closes most of the gap to the hybrid chain.
\item phase2\_rf\_fair\_changed: Fair RF budgeting rescues a branch that earlier looked much weaker.
\item phase3b\_raw\_only\_changed: Rollout-aligned retraining recovers a strong raw chain instead of chance collapse.
\item phase1\_patch\_overlap\_changed: Patchwise overlap materially improves over the pointwise vector baseline.
\item phase1\_patch\_depth\_changed: Depth-aware schedules complete cleanly and broadly support the small-context story.
\item phase4\_r20c10: The winner recipe survives the ResNet-20 / CIFAR-10 transfer grid.
\item phase4\_seeds: Seed variation is small enough that the qualitative story survives.
\item phase4\_taxonomy: Sparse is the best transition carrier among the tested alternatives.
\item phase4\_r110: The winner recipe transfers to a deeper CNN and preserves the Q3 bottleneck pattern.
\item phase4\_vit\_transfer: The ViT transfer artifact achieves high prediction agreement with low KL at 5 percent retained.
\item phase4\_vit\_sweep: The completed ViT-small sweep shows that the transformer-side result is stable across multiple blocks, sparsity budgets, and seeds, with a best stable setting at block 10 and 12 percent retained.
\item phase4\_vit\_base: The completed ViT-base sweep shows a much stronger stable best-case regime than ViT-small, but also surfaces substantial seed-sensitive fragility outside the late block-10 regime.
\item tier1\_intervention: Interventions show strong importance and locality concentration, though they inherit upstream SAE validity risks.
\item tier2\_feature\_audit: Dead-feature and duplicate-feature rates are low in the audited sections.
\end{itemize}
\section{What still should be run}
\begin{itemize}
\item A validation-first rerun of the CIFAR families that removes repeated test-set evaluation from training loops, if we later want a separate external-benchmarking layer on top of the mechanistic claim.
\item A regenerated Family I chain artifact that stores raw, re-mask, and re-ground chain variants as standalone verified outputs instead of embedded report values.
\item If we later want untouched external evaluation numbers, reserve a fresh held-out split or a fully untouched test set for final scoring.
\end{itemize}
\section{Next steps}
\begin{enumerate}
\item If we want an additional external-benchmarking layer, remove test-set evaluation from the SAE, transition, and chain training loops, and rerun the headline CIFAR grids against a clean validation split.
\item Keep the FieldSAE winner and the re-grounded chain path as the main line; do not spend more time on coarse disjoint patches or simple larger Q3-to-Q4 predictors.
\item Use the ViT path as the cleanest transfer scaffold, because its surviving artifact and completed sweep do not show the same split-discipline problem as the CIFAR training families.
\item Prioritize feature-visualization and seed-universality follow-ups on the stable ViT winners, and separately diagnose the realized-versus-requested sparsity drift behind the fragile ViT-base cells before broadening the claim to still larger models or datasets.
\end{enumerate}
\section{Appendix}
Primary leakage evidence: src/train\_sae.py:77-78,110-117; src/train\_sae\_changed.py:143-148,196-205; src/train\_transition.py:71-72,87-92; src/train\_chain.py:61-63,101-105.


Split-hygiene evidence: src/data.py:18-30; src/cache\_activations.py:37-45,57-60; src/vit\_sae.py:65-72,79-81,96-112.


\end{document}
